\section{Problem, scope, and epistemic state}

\subsection{The failure mode}

A research process can make progress while its own framing becomes the dominant source of error. More search does not repair a missing representation. A larger tree does not by itself reveal that the wrong parent discipline has been treated as irrelevant. More evidence does not justify rewriting a question when the actual defect is merely an execution failure. Paper I studies the resulting governance problem: preserving an inspectable epistemic state while deciding when evidence is strong enough to authorize a change in the scientific formulation itself.

For a scoped problem, ORION represents the evolving research state as
\[
X_t=(K_t,W_t,M_t),
\]
where $K_t$ is the provenance-preserving state of knowledge about the target object, $W_t$ is the current model of which domains, representations, mechanisms, and search routes may be relevant, and $M_t$ is the governed research method. The three coordinates are separated because they require different evidence and different authority to change. This engineering separation is in the spirit of Parnas's information-hiding principle: change-prone design decisions should be isolated behind explicit interfaces rather than entangled implicitly \cite{parnas1972}. The revision policy for $K_t$ likewise borrows the minimal-change distinction between expansion, contraction, and revision from the AGM belief-revision tradition \cite{agm1985}; Paper I does not claim that its operational update rules satisfy the full AGM postulates unless that correspondence is established separately.

A new source may update $K_t$ without changing $W_t$. A previously unrecognized parent discipline or representation may require a change to $W_t$ while leaving the method intact. A persistent, well-attributed failure can become evidence for a challenger to $M_t$, but method revision is governed separately by Self-ORION and is not a unilateral Paper-I operation.

\subsection{Tested contribution}

The final programme isolates one dominant contribution with four coupled requirements. First, $K/W/M$ are explicit co-evolving state rather than hidden conversational context. Second, discoveries and failures target typed formulation coordinates, so diagnosed responsibility constrains which high-level mutation is even admissible. Third, high-level change is gated by lower-level exclusion, same-task counterfactual necessity, and protected-invariant preservation rather than by generic ``reflect and retry.'' Fourth, an accepted mutation stales exactly the dependent closure identified by its bound impact rather than allowing prior success certificates to survive silently.

The empirical studies test this composition directly. On hidden formulation shifts, the protected necessity contract separates successful reframing from indiscriminate reframing; on evidence/execution controls, the same contract tests whether ORION can refrain from changing the formulation when a lower-level repair is sufficient. The P1-X successor then generalizes the same idea from one hidden-shift mechanism to heterogeneous revision-responsibility contracts.

The contribution is therefore not ownership of iteration, tree search, recursive audit, dependency repair, or generic mutation authorization. Those mechanisms are absorbed as donor substrate. The scientific claim is that coupling responsibility, counterfactual necessity, protected preservation, and dependency impact creates a stronger decision contract for high-level scientific reformulation, and that this contract produces large protected gains on the prospectively frozen families evaluated here.

\subsection{Five-paper boundary}

Paper I owns the reconstruction problem only. Open-world discovery and stopping are evaluated in Paper II; cross-domain absorption and the global knowledge portrait in Paper III; authority promotion and protected verification in Paper IV; and method evolution in Paper V. Paper I uses those components only as interfaces required to express reconstruction. This separation is important because an umbrella architecture description would otherwise allow the same mechanism or empirical result to support several nominally different papers.
